% Transcription diplomatique française de EGA IV, quatrième partie, pp. 341--343.
\clearpage
\oldpage[IV]{341}
\phantomsection
\addcontentsline{toc}{section}{Table des matières}
\section*{TABLE DES MATIÈRES}

\begingroup
\small
\newcommand{\EGAIVcontentsline}[3]{%
  \textbf{#1}\quad #2\dotfill #3\par}

\begin{center}\textbf{CHAPITRE IV. --- Étude locale des schémas et des morphismes de schémas (fin)}\end{center}

\EGAIVcontentsline{\S16}{Invariants différentiels. Morphismes différentiellement lisses}{5}
\EGAIVcontentsline{{\texorpdfstring{\protect\hyperlink{ega4.c04.subsec.16-1}{16.1}}{16.1}}}{Invariants normaux d'une immersion}{5}
\EGAIVcontentsline{{\texorpdfstring{\protect\hyperlink{ega4.c04.subsec.16-2}{16.2}}{16.2}}}{Propriétés fonctorielles des invariants normaux d'une immersion}{9}
\EGAIVcontentsline{{\texorpdfstring{\protect\hyperlink{ega4.c04.subsec.16-3}{16.3}}{16.3}}}{Invariants différentiels fondamentaux d'un morphisme de préschémas}{14}
\EGAIVcontentsline{{\texorpdfstring{\protect\hyperlink{ega4.c04.subsec.16-4}{16.4}}{16.4}}}{Propriétés fonctorielles des invariants différentiels}{16}
\EGAIVcontentsline{{\texorpdfstring{\protect\hyperlink{ega4.c04.subsec.16-5}{16.5}}{16.5}}}{Faisceaux et fibrés tangents relatifs; dérivations}{27}
\EGAIVcontentsline{{\texorpdfstring{\protect\hyperlink{ega4.c04.subsec.16-6}{16.6}}{16.6}}}{Faisceaux de $p$-différentielles et différentielle extérieure}{34}
\EGAIVcontentsline{{\texorpdfstring{\protect\hyperlink{ega4.c04.subsec.16-7}{16.7}}{16.7}}}{Les $\sh{P}^{n}_{X/S}(\sh{F})$}{36}
\EGAIVcontentsline{{\texorpdfstring{\protect\hyperlink{ega4.c04.subsec.16-8}{16.8}}{16.8}}}{Opérateurs différentiels}{39}
\EGAIVcontentsline{{\texorpdfstring{\protect\hyperlink{ega4.c04.subsec.16-9}{16.9}}{16.9}}}{Immersions régulières et quasi-régulières}{46}
\EGAIVcontentsline{{\texorpdfstring{\protect\hyperlink{ega4.c04.subsec.16-10}{16.10}}{16.10}}}{Morphismes différentiellement lisses}{51}
\EGAIVcontentsline{{\texorpdfstring{\protect\hyperlink{ega4.c04.subsec.16-11}{16.11}}{16.11}}}{Opérateurs différentiels sur un $S$-préschéma différentiellement lisse}{53}
\EGAIVcontentsline{{\texorpdfstring{\protect\hyperlink{ega4.c04.subsec.16-12}{16.12}}{16.12}}}{Cas de la caractéristique nulle : critère jacobien pour les morphismes différentiellement lisses}{55}

\EGAIVcontentsline{\S17}{Morphismes lisses, morphismes non ramifiés, morphismes étales}{56}
\EGAIVcontentsline{{\texorpdfstring{\protect\hyperlink{ega4.c04.subsec.17-1}{17.1}}{17.1}}}{Morphismes formellement lisses, morphismes formellement non ramifiés, morphismes formellement étales}{56}
\EGAIVcontentsline{{\texorpdfstring{\protect\hyperlink{ega4.c04.subsec.17-2}{17.2}}{17.2}}}{Propriétés différentielles générales}{59}
\EGAIVcontentsline{{\texorpdfstring{\protect\hyperlink{ega4.c04.subsec.17-3}{17.3}}{17.3}}}{Morphismes lisses, morphismes non ramifiés, morphismes étales}{61}
\EGAIVcontentsline{{\texorpdfstring{\protect\hyperlink{ega4.c04.subsec.17-4}{17.4}}{17.4}}}{Caractérisations des morphismes non ramifiés}{65}
\EGAIVcontentsline{{\texorpdfstring{\protect\hyperlink{ega4.c04.subsec.17-5}{17.5}}{17.5}}}{Caractérisations des morphismes lisses}{67}
\EGAIVcontentsline{{\texorpdfstring{\protect\hyperlink{ega4.c04.subsec.17-6}{17.6}}{17.6}}}{Caractérisations des morphismes étales}{70}
\EGAIVcontentsline{{\texorpdfstring{\protect\hyperlink{ega4.c04.subsec.17-7}{17.7}}{17.7}}}{Propriétés de descente et de passage à la limite}{72}
\EGAIVcontentsline{{\texorpdfstring{\protect\hyperlink{ega4.c04.subsec.17-8}{17.8}}{17.8}}}{Critères de lissité et de non ramification par fibres}{79}
\EGAIVcontentsline{{\texorpdfstring{\protect\hyperlink{ega4.c04.subsec.17-9}{17.9}}{17.9}}}{Morphismes étales et immersions ouvertes}{79}
\EGAIVcontentsline{{\texorpdfstring{\protect\hyperlink{ega4.c04.subsec.17-10}{17.10}}{17.10}}}{Dimension relative d'un préschéma lisse sur un autre}{81}
\EGAIVcontentsline{{\texorpdfstring{\protect\hyperlink{ega4.c04.subsec.17-11}{17.11}}{17.11}}}{Morphismes lisses de préschémas lisses}{82}
\EGAIVcontentsline{{\texorpdfstring{\protect\hyperlink{ega4.c04.subsec.17-12}{17.12}}{17.12}}}{Sous-préschémas lisses d'un préschéma lisse. Morphismes lisses et morphismes différentiellement lisses}{85}

\oldpage[IV]{342}
\EGAIVcontentsline{{\texorpdfstring{\protect\hyperlink{ega4.c04.subsec.17-13}{17.13}}{17.13}}}{Morphismes transversaux}{89}
\EGAIVcontentsline{{\texorpdfstring{\protect\hyperlink{ega4.c04.subsec.17-14}{17.14}}{17.14}}}{Caractérisations locales et infinitésimales des morphismes lisses, des morphismes non ramifiés et des morphismes étales}{98}
\EGAIVcontentsline{{\texorpdfstring{\protect\hyperlink{ega4.c04.subsec.17-15}{17.15}}{17.15}}}{Cas des préschémas sur un corps de base}{99}
\EGAIVcontentsline{{\texorpdfstring{\protect\hyperlink{ega4.c04.subsec.17-16}{17.16}}{17.16}}}{Quasi-sections de morphismes plats ou lisses}{105}

\EGAIVcontentsline{\S18}{Compléments sur les morphismes étales. Anneaux locaux henséliens et anneaux strictement locaux}{109}
\EGAIVcontentsline{{\texorpdfstring{\protect\hyperlink{ega4.c04.subsec.18-1}{18.1}}{18.1}}}{Une équivalence remarquable de catégories}{109}
\EGAIVcontentsline{{\texorpdfstring{\protect\hyperlink{ega4.c04.subsec.18-2}{18.2}}{18.2}}}{Revêtements étales}{111}
\EGAIVcontentsline{{\texorpdfstring{\protect\hyperlink{ega4.c04.subsec.18-3}{18.3}}{18.3}}}{Algèbres finies et étales}{114}
\EGAIVcontentsline{{\texorpdfstring{\protect\hyperlink{ega4.c04.subsec.18-4}{18.4}}{18.4}}}{Structure locale des morphismes non ramifiés et des morphismes étales}{118}
\EGAIVcontentsline{{\texorpdfstring{\protect\hyperlink{ega4.c04.subsec.18-5}{18.5}}{18.5}}}{Anneaux locaux henséliens}{125}
\EGAIVcontentsline{{\texorpdfstring{\protect\hyperlink{ega4.c04.subsec.18-6}{18.6}}{18.6}}}{Hensélisation}{135}
\EGAIVcontentsline{{\texorpdfstring{\protect\hyperlink{ega4.c04.subsec.18-7}{18.7}}{18.7}}}{Hensélisation et anneaux excellents}{142}
\EGAIVcontentsline{{\texorpdfstring{\protect\hyperlink{ega4.c04.subsec.18-8}{18.8}}{18.8}}}{Anneaux strictement locaux et hensélisation stricte}{144}
\EGAIVcontentsline{{\texorpdfstring{\protect\hyperlink{ega4.c04.subsec.18-9}{18.9}}{18.9}}}{Fibres formelles des anneaux noethériens henséliens}{150}
\EGAIVcontentsline{{\texorpdfstring{\protect\hyperlink{ega4.c04.subsec.18-10}{18.10}}{18.10}}}{Préschémas étales sur un préschéma géométriquement unibranche ou normal}{157}
\EGAIVcontentsline{{\texorpdfstring{\protect\hyperlink{ega4.c04.subsec.18-11}{18.11}}{18.11}}}{Application aux algèbres locales noethériennes complètes sur un corps}{169}
\EGAIVcontentsline{{\texorpdfstring{\protect\hyperlink{ega4.c04.subsec.18-12}{18.12}}{18.12}}}{Applications de la localisation étale aux morphismes quasi-finis (généralisations de résultats antérieurs)}{181}

\EGAIVcontentsline{\S19}{Immersions régulières et platitude normale}{185}
\EGAIVcontentsline{{\texorpdfstring{\protect\hyperlink{ega4.c04.subsec.19-1}{19.1}}{19.1}}}{Propriétés des immersions régulières}{185}
\EGAIVcontentsline{{\texorpdfstring{\protect\hyperlink{ega4.c04.subsec.19-2}{19.2}}{19.2}}}{Immersions transversalement régulières}{190}
\EGAIVcontentsline{{\texorpdfstring{\protect\hyperlink{ega4.c04.subsec.19-3}{19.3}}{19.3}}}{Intersections complètes relatives (cas plat)}{194}
\EGAIVcontentsline{{\texorpdfstring{\protect\hyperlink{ega4.c04.subsec.19-4}{19.4}}{19.4}}}{Application : critères de régularité et de lissité pour les préschémas éclatés}{198}
\EGAIVcontentsline{{\texorpdfstring{\protect\hyperlink{ega4.c04.subsec.19-5}{19.5}}{19.5}}}{Critères de $M$-régularité}{204}
\EGAIVcontentsline{{\texorpdfstring{\protect\hyperlink{ega4.c04.subsec.19-6}{19.6}}{19.6}}}{Suites régulières relativement à un module filtré quotient}{209}
\EGAIVcontentsline{{\texorpdfstring{\protect\hyperlink{ega4.c04.subsec.19-7}{19.7}}{19.7}}}{Critère de platitude normale de Hironaka}{212}
\EGAIVcontentsline{{\texorpdfstring{\protect\hyperlink{ega4.c04.subsec.19-8}{19.8}}{19.8}}}{Propriétés de passage à la limite projective}{219}
\EGAIVcontentsline{{\texorpdfstring{\protect\hyperlink{ega4.c04.subsec.19-9}{19.9}}{19.9}}}{Suites $\sh{F}$-régulières et profondeur}{222}

\EGAIVcontentsline{\S20}{Fonctions méromorphes et pseudo-morphismes}{223}
\EGAIVcontentsline{{\texorpdfstring{\protect\hyperlink{ega4.c04.subsec.20-0}{20.0}}{20.0}}}{Introduction}{223}
\EGAIVcontentsline{{\texorpdfstring{\protect\hyperlink{ega4.c04.subsec.20-1}{20.1}}{20.1}}}{Fonctions méromorphes}{224}
\EGAIVcontentsline{{\texorpdfstring{\protect\hyperlink{ega4.c04.subsec.20-2}{20.2}}{20.2}}}{Pseudo-morphismes et pseudo-fonctions}{231}

\oldpage[IV]{343}
\EGAIVcontentsline{{\texorpdfstring{\protect\hyperlink{ega4.c04.subsec.20-3}{20.3}}{20.3}}}{Composition des pseudo-morphismes}{237}
\EGAIVcontentsline{{\texorpdfstring{\protect\hyperlink{ega4.c04.subsec.20-4}{20.4}}{20.4}}}{Propriétés des domaines de définition des fonctions rationnelles}{244}
\EGAIVcontentsline{{\texorpdfstring{\protect\hyperlink{ega4.c04.subsec.20-5}{20.5}}{20.5}}}{Pseudo-morphismes relatifs}{249}
\EGAIVcontentsline{{\texorpdfstring{\protect\hyperlink{ega4.c04.subsec.20-6}{20.6}}{20.6}}}{Fonctions méromorphes relatives}{252}

\EGAIVcontentsline{\S21}{Diviseurs}{255}
\EGAIVcontentsline{{\texorpdfstring{\protect\hyperlink{ega4.c04.subsec.21-1}{21.1}}{21.1}}}{Diviseurs sur un espace annelé}{255}
\EGAIVcontentsline{{\texorpdfstring{\protect\hyperlink{ega4.c04.subsec.21-2}{21.2}}{21.2}}}{Diviseurs et Idéaux fractionnaires inversibles}{258}
\EGAIVcontentsline{{\texorpdfstring{\protect\hyperlink{ega4.c04.subsec.21-3}{21.3}}{21.3}}}{Équivalence linéaire des diviseurs}{263}
\EGAIVcontentsline{{\texorpdfstring{\protect\hyperlink{ega4.c04.subsec.21-4}{21.4}}{21.4}}}{Images réciproques de diviseurs}{265}
\EGAIVcontentsline{{\texorpdfstring{\protect\hyperlink{ega4.c04.subsec.21-5}{21.5}}{21.5}}}{Images directes de diviseurs}{267}
\EGAIVcontentsline{{\texorpdfstring{\protect\hyperlink{ega4.c04.subsec.21-6}{21.6}}{21.6}}}{Cycle 1-codimensionnel associé à un diviseur}{270}
\EGAIVcontentsline{{\texorpdfstring{\protect\hyperlink{ega4.c04.subsec.21-7}{21.7}}{21.7}}}{Interprétation des cycles positifs 1-codimensionnels en termes de sous-préschémas}{277}
\EGAIVcontentsline{{\texorpdfstring{\protect\hyperlink{ega4.c04.subsec.21-8}{21.8}}{21.8}}}{Diviseurs et normalisation}{280}
\EGAIVcontentsline{{\texorpdfstring{\protect\hyperlink{ega4.c04.subsec.21-9}{21.9}}{21.9}}}{Diviseurs sur les préschémas de dimension 1}{284}
\EGAIVcontentsline{{\texorpdfstring{\protect\hyperlink{ega4.c04.subsec.21-10}{21.10}}{21.10}}}{Images réciproques et images directes de cycles 1-codimensionnels}{289}
\EGAIVcontentsline{{\texorpdfstring{\protect\hyperlink{ega4.c04.subsec.21-11}{21.11}}{21.11}}}{Factorialité des anneaux réguliers}{302}
\EGAIVcontentsline{{\texorpdfstring{\protect\hyperlink{ega4.c04.subsec.21-12}{21.12}}{21.12}}}{Le théorème de pureté de Van der Waerden pour l'ensemble de ramification d'un morphisme birationnel}{304}
\EGAIVcontentsline{{\texorpdfstring{\protect\hyperlink{ega4.c04.subsec.21-13}{21.13}}{21.13}}}{Couples parafactoriels. Anneaux locaux parafactoriels}{313}
\EGAIVcontentsline{{\texorpdfstring{\protect\hyperlink{ega4.c04.subsec.21-14}{21.14}}{21.14}}}{Le théorème de Ramanujam--Samuel}{323}
\EGAIVcontentsline{{\texorpdfstring{\protect\hyperlink{ega4.c04.subsec.21-15}{21.15}}{21.15}}}{Diviseurs relatifs}{329}
\EGAIVcontentsline{}{BIBLIOGRAPHIE}{333}
\EGAIVcontentsline{}{INDEX DES NOTATIONS}{334}
\EGAIVcontentsline{}{INDEX TERMINOLOGIQUE}{336}

\endgroup
